\subsubsection{Strategic Trading Models}

\begin{taxonomy} \hlt{Informed Trader Style Segmentation}
\begin{enumerate}[label=\roman*.]
\setlength{\itemsep}{0pt}
\item \hlt{value}: estimate total fundamental value from broad information;
\item \hlt{news}: estimate value changes from recent information;
\item \hlt{information-oriented technical}: exploit predictable errors by other traders;
\item \hlt{arbitrage}: trade relative mispricing across instruments with common valuation factors.
\end{enumerate}
\end{taxonomy}

\begin{taxonomy} \hlt{Informed Trading Speed Segmentation}
\begin{enumerate}[label=\roman*.]
\setlength{\itemsep}{0pt}
\item \hlt{aggressive}: trade quickly when information is short-lived or competitors know it;
\item \hlt{stealth}: trade slowly when impact and information leakage dominate urgency;
\item \hlt{inactive}: avoid trading when the price already reflects the information.
\end{enumerate}
\end{taxonomy}

\begin{principle} \hlt{Orthogonal Information Advantage}\\
The best informed traders are not only accurate; their errors are different from other traders' errors. Orthogonal information is more valuable than crowded information with the same average accuracy.
\end{principle}

\begin{caveat} \hlt{Informed Trading Paradox}\\
Perfectly informative prices eliminate informed profits; zero informed profits stop prices from becoming informative. Markets persist between these limits because information is costly, temporary, and unevenly distributed.
\end{caveat}
